\chapter{Introduction to Scientific Workflows}
\label{ch:introduction}

\epigraph{If these steps remain implicit in shell history, personal notebooks,
or site-specific scripts, the experiment is difficult to audit. Workflow
systems make the structure explicit.}{\textit{DAGonStar documentation,
Introduction to Scientific Workflows}}

\section{Why Workflows Matter in Computational Science}

Modern computational science rarely consists of a single program executed
once. A realistic environmental-science experiment, for instance, combines
many heterogeneous activities:

\begin{itemize}[leftmargin=1.4em]
  \item preprocessing observational data;
  \item running models or simulations;
  \item staging intermediate files between machines;
  \item executing tools on HPC clusters or cloud instances;
  \item post-processing model products;
  \item publishing scientific outputs;
  \item preserving sufficient context to reproduce a result.
\end{itemize}

When these steps are encoded only in shell history, ad hoc notebooks, or
site-specific scripts passed from one student to the next, the resulting
experiment is opaque: it cannot easily be audited, repeated by a different
researcher, or ported to a different machine. A \emph{workflow system} exists
precisely to make this implicit structure explicit, executable, and
inspectable.

DAGonStar approaches the problem with a deliberately small abstraction: a
\textbf{workflow is a directed acyclic graph (DAG) of tasks}, and each
\textbf{task} is a Python object that executes a command in a particular
environment. The remainder of this chapter develops this abstraction
formally and explains the scientific benefits it is designed to deliver.

\subsection{The Scientific Benefits of Explicit Workflows}

The DAGonStar documentation identifies five principal benefits that follow
from making computational structure explicit:

\begin{description}[leftmargin=1.6em,style=nextline]
  \item[Reproducibility] the sequence of computations is encoded as source
    code, not as institutional memory.
  \item[Traceability] inputs, outputs, and task dependencies are visible as
    first-class graph structure rather than buried in script side effects.
  \item[Portability] the same workflow structure can target different
    execution backends --- a laptop today, an HPC cluster tomorrow --- without
    rewriting the scientific logic.
  \item[Parallelism] independent tasks can run concurrently when resources
    permit, because the dependency graph makes independence explicit.
  \item[Recoverability] checkpointing can avoid recomputing finished work
    after a partial failure, which matters enormously for expensive
    simulations.
\end{description}

\begin{practicebox}
Even apparently small automation gains compound at scale. A workflow that
takes ten extra minutes to encode explicitly, instead of being run by hand
from a terminal, pays for itself the first time it must be rerun --- with a
different parameter, on a different machine, or six months after the original
researcher has left the project.
\end{practicebox}

\section{Directed Acyclic Graphs}
\label{sec:dag-theory}

\begin{definitionbox}{Directed Acyclic Graph (DAG)}
A \emph{directed acyclic graph} $G = (V, E)$ consists of a finite set of
\textbf{nodes} $V$ and a set of directed \textbf{edges} $E \subseteq V \times
V$, such that there is no sequence of edges
$v_1 \to v_2 \to \dots \to v_k \to v_1$ that returns to its starting node. In
scientific workflow terminology:
\begin{itemize}
  \item a \textbf{node} is a computational \emph{task};
  \item an \textbf{edge} $(u,v) \in E$ means that task $v$ \emph{depends on}
    task $u$, i.e.\ $u$ must complete before $v$ may begin;
  \item \textbf{acyclic} means the graph contains no circular dependency.
\end{itemize}
\end{definitionbox}

The acyclicity constraint is not a mathematical nicety: it is the precise
condition under which a valid execution order (a \emph{topological sort})
exists. If a cycle were permitted, no task in the cycle could be scheduled
first without violating some dependency in the cycle, and the workflow would
be unrealizable.

\begin{figure}[h]
\centering
\begin{tikzpicture}[
  node distance=16mm and 22mm,
  task/.style={rectangle, rounded corners, draw=dagonblue, very thick,
    fill=blue!4, minimum width=20mm, minimum height=9mm,
    font=\sffamily\bfseries},
  arr/.style={-{Stealth[length=2.6mm]}, thick, dagonblue}
]
  \node[task] (A) {Tokio};
  \node[task, right=of A] (D) {Mosco};
  \node[task, above=of D] (B) {Berlin};
  \node[task, below=of D] (C) {Nairobi};
  \draw[arr] (A) -- (B);
  \draw[arr] (A) -- (C);
  \draw[arr] (B) -- (D);
  \draw[arr] (C) -- (D);
\end{tikzpicture}
\caption{A four-task DAG corresponding to the engine's canonical
``Batch Task Flow'' example. \texttt{Tokio} has no predecessors and runs
first; \texttt{Berlin} and \texttt{Nairobi} depend only on \texttt{Tokio} and
may run concurrently; \texttt{Mosco} depends on both and runs last.}
\label{fig:basic-dag}
\end{figure}

\Cref{fig:basic-dag} illustrates the canonical example used throughout the
DAGonStar project, expressed as four \texttt{BATCH} tasks: \texttt{Tokio}
writes a file consumed implicitly by \texttt{Mosco}, while \texttt{Berlin} and
\texttt{Nairobi} are independent siblings that both depend only on
\texttt{Tokio}. The graph immediately exposes the available parallelism: a
scheduler is free to run \texttt{Berlin} and \texttt{Nairobi} concurrently
once \texttt{Tokio} finishes.

DAGonStar validates the workflow graph before execution, either implicitly
when dependencies are constructed through \texttt{Workflow.make\_dependencies()},
or explicitly via \texttt{Workflow.Validate\_WF()}. A cycle is treated as a
\emph{modeling error}: the engine refuses to run a workflow whose dependency
graph is not a true DAG.

\section{Task-Oriented and Data-Oriented Workflows}
\label{sec:task-vs-data}

DAGonStar supports two complementary ways of describing a workflow, and the
distinction between them is the single most important conceptual choice a
workflow author makes.

\subsection{Task-Oriented Workflows}

In a \textbf{task-oriented workflow}, dependencies are stated explicitly by
the author, independent of any file the tasks happen to exchange:

\begin{lstlisting}[caption={Explicit dependency declaration.}]
task_b.add_dependency_to(task_a)
\end{lstlisting}

This statement reads: ``\texttt{task\_b} must wait for \texttt{task\_a}.''
Explicit dependencies are appropriate when the relationship between two tasks
is \emph{conceptual or procedural} rather than a literal file consumption ---
for example, ``validate parameters before launching the model.''

\subsection{Data-Oriented Workflows}

In a \textbf{data-oriented workflow}, dependencies are \emph{inferred} from
the command text using the \texttt{workflow://} schema (developed fully in
\Cref{ch:schema}):

\begin{lstlisting}[caption={Implicit, data-driven dependency.}]
task_b = DagonTask(TaskType.BATCH, "B", "cat workflow:///A/output/result.txt")
\end{lstlisting}

When \texttt{Workflow.make\_dependencies()} is called, DAGonStar scans every
task's command string, finds the reference to task \texttt{A}, records the
dependency $A \to B$, and later stages the referenced file into \texttt{B}'s
working directory before \texttt{B} executes.

\begin{practicebox}
The data-oriented style is especially natural for scientific pipelines,
because dependency edges then follow the \emph{data products} being consumed
--- which is, after all, how a scientist already thinks about a processing
chain: ``the post-processing step needs the simulation's NetCDF output,''
not ``task three must follow task two for an unstated reason.''
\end{practicebox}

Both styles may, and routinely do, coexist in the same workflow: explicit
edges encode procedural ordering, while \texttt{workflow://} edges encode
genuine data lineage. \Cref{ch:lessons-01-02} develops both patterns with
complete, runnable examples.

\section{DAGonStar's Computational Model}
\label{sec:computational-model}

The current implementation centers on a small number of cooperating
abstractions, summarized here and developed individually across the
remainder of Part~I and Part~II:

\begin{description}[leftmargin=2.2em,style=nextline]
  \item[\texttt{Workflow}] the container and scheduler for a set of tasks;
    owns dependency construction, execution, checkpoint state, and
    configuration.
  \item[\texttt{DagonTask}] a \emph{factory} that constructs a concrete task
    implementation from a \texttt{TaskType} value.
  \item[\texttt{Task}] the base class providing threaded task execution,
    dependency bookkeeping, and \texttt{workflow://} resolution.
  \item[\texttt{Batch}] a local shell task --- the simplest concrete
    execution backend.
  \item[\texttt{Slurm}] a task submitted through \texttt{sbatch} to an HPC
    job scheduler.
  \item[\texttt{RemoteTask}] a task executed on a remote host through SSH.
  \item[\texttt{DockerTask}] a task executed inside a Docker container.
  \item[\texttt{CloudTask}] a task bound to an Apache Libcloud-managed cloud
    instance.
  \item[\texttt{Stager}] the component responsible for generating or
    invoking the data-movement operation between two dependent tasks.
\end{description}

Operationally, the workflow engine runs tasks as Python \texttt{threading.Thread}
objects. Each task thread waits for its predecessors, prepares its scratch
directory, resolves data references, executes its generated launcher script,
updates its status, and finally releases the tasks that depend on it. This
execution model is examined in full mechanical detail in
\Cref{ch:architecture}.

\section{The Role of Scratch Directories}

Scientific workflows frequently produce many intermediate files. DAGonStar
therefore gives \emph{every} task its own working directory, called a
\textbf{scratch directory}. The base path for all scratch directories is set
once, in configuration:

\begin{lstlisting}[style=ini,caption={Minimal scratch directory configuration.}]
[batch]
scratch_dir_base=/tmp/
\end{lstlisting}

If a task is not given an explicit \texttt{working\_dir}, DAGonStar
constructs one automatically from the workflow's scratch base and a
timestamped task name (for example, \texttt{/tmp/1710000000000-prepare}).
Every task additionally receives an internal \texttt{.dagon/} directory
containing its generated launcher script and execution metadata --- a
mechanism examined closely in \Cref{ch:lessons-03-04}.

\section{Execution Backends at a Glance}

DAGonStar's public task types are enumerated in \texttt{dagon.task.TaskType}.
There are sixteen in DAGonStar~3.0.0 (\texttt{examples/03-architecture}); the
main ones are:

\begin{center}
\begin{tabular}{@{}ll@{}}
\toprule
\textbf{\texttt{TaskType}} & \textbf{Execution environment} \\
\midrule
\texttt{BATCH}      & Local (or remote, via SSH keyword arguments) shell command \\
\texttt{SLURM}       & HPC job-scheduler submission via \texttt{sbatch} \\
\texttt{CLOUD}       & Apache Libcloud-managed virtual machine \\
\texttt{DOCKER}      & Containerized execution via the Docker SDK \\
\texttt{LLM}         & OpenAI-compatible Chat Completions request \\
\texttt{NATIVE}      & Importable Python \texttt{module:function} call \\
\texttt{WEB}         & Arbitrary HTTP/HTTPS request \\
\texttt{CHECKPOINT}  & Explicit checkpoint-preservation task \\
\texttt{META}        & A workflow of another engine, run unmodified (Appendix~\ref{app:meta-tasks}) \\
\texttt{VOLUME}      & Creates or destroys a shared ad-hoc volume (Appendix~\ref{app:volumes}) \\
\bottomrule
\end{tabular}
\end{center}

Remote variants of several task types are selected not by a different
\texttt{TaskType} value, but by \emph{keyword arguments}: for example, a
\texttt{BATCH} task constructed with \texttt{ip=...} becomes a remote batch
task transparently. This design keeps workflow code compact while still
exposing several execution environments through one uniform factory,
\texttt{DagonTask}. It also means that several integrations are
\emph{optional}: using Docker, cloud, Globus, or remote SSH execution
requires the corresponding site-specific service or credential, while the
core local-batch engine has no such external dependency.

\begin{practicebox}
\texttt{LLM} tasks --- the engine's principal AI integration point --- call
configured OpenAI-compatible Chat Completions providers and can consume local
UTF-8 producer files through the same \texttt{workflow://} schema used by
ordinary file-based tasks. This is the technical mechanism that allows an AI
inference step to sit inside the same DAG as a numerical simulation, with the
same dependency, staging, and provenance guarantees. \Cref{ch:llm-tasks} is
devoted entirely to this mechanism.
\end{practicebox}

\section{Workflows in the Landscape of Computational Science}

It is useful, before going further, to situate the workflow-orchestration
problem within the broader history of computational science tooling, so
that the design choices examined in later chapters can be understood not as
arbitrary engineering preferences but as responses to a recurring,
well-documented set of pressures.

Three historically distinct strategies have been used to compose
computational steps into a pipeline. The first, and oldest, is the
\emph{shell script}: a linear sequence of commands, executed in source-code
order, with control flow expressed through exit codes and conditional
branches. Shell scripts are easy to write and require no additional
software, but they conflate two logically separate concerns --- the
\emph{order} in which steps must run and the \emph{mechanism} by which they
run --- into a single, brittle, line-by-line narrative. A shell script
cannot, without substantial additional engineering, express that two steps
are mutually independent and may run concurrently, nor can it express, in a
form a second program can inspect, which step produced which file.

The second strategy is the \emph{build system}, exemplified historically by
\texttt{make} and its descendants. Build systems were designed to compile
software, not to run scientific experiments, but they introduced an idea
central to this book: a target (a file to be produced) depends on a set of
prerequisites (other files), and a recipe (a command) explains how to
produce the target from the prerequisites once they exist. This is, in
essence, a restricted form of the data-oriented dependency model developed
in \Cref{sec:task-vs-data}. Build systems are excellent at avoiding
redundant recomputation --- a target whose prerequisites have not changed is
not rebuilt --- but they are typically confined to a single machine and a
single execution environment, and their dependency language is tied tightly
to file modification timestamps rather than to a richer notion of
provenance.

The third strategy, and the one this book is concerned with, is the
\emph{dedicated scientific workflow engine}. Such engines generalize the
build-system insight --- dependencies as a directed graph over named
artifacts --- while relaxing its constraints: a task may run on a different
machine than its consumer, may be backed by a container, a job scheduler, a
remote web service, or an AI inference call, and the engine is expected to
manage the data movement, environment provisioning, and failure recovery
that a single-machine build tool can ignore. DAGonStar belongs to this third
family, and the remainder of this book can be read as an extended case
study in how one engine in this family resolves the tension between
\emph{generality} (supporting arbitrary execution backends) and
\emph{simplicity} (keeping the authoring surface close to ordinary Python).

\begin{practicebox}
A useful heuristic for deciding whether a computational pipeline has outgrown
a shell script and needs a workflow engine: if you find yourself writing
comments in a shell script that explain \emph{why} step three must follow
step one, or maintaining a parallel README describing which script produced
which file, the dependency structure you are documenting informally is
exactly the structure a DAG-based engine would make executable.
\end{practicebox}

\section{A Formal Note on Topological Ordering}
\label{sec:topological-ordering}

\Cref{sec:dag-theory} asserted, without proof, that acyclicity is precisely
the condition under which a valid execution order exists. It is worth
making this claim mathematically precise, both because it explains
\emph{why} \texttt{Validate\_WF()} is sufficient to guarantee an executable
workflow, and because it grounds the engine's threaded execution model
(\Cref{sec:execution-lifecycle}) in a standard result of graph theory.

\begin{definitionbox}{Topological Order}
A \emph{topological order} of a directed graph $G=(V,E)$ is a linear
ordering $v_1, v_2, \dots, v_n$ of $V$ such that, for every edge
$(v_i, v_j) \in E$, $i < j$. Equivalently, every node appears in the
ordering only after all nodes it depends on.
\end{definitionbox}

A standard theorem of graph theory states that a directed graph admits a
topological order \emph{if and only if} it is acyclic. The forward direction
is immediate by contradiction: if a cycle $v_1 \to v_2 \to \dots \to v_k \to
v_1$ existed, no linear ordering could place every predecessor before its
successor along the entire cycle, since the cycle returns to its starting
node. The converse can be shown constructively, by repeated removal of
\emph{source} nodes --- nodes with no remaining incoming edges --- from the
graph: at every step, an acyclic graph with at least one node must contain
at least one source (otherwise, following incoming edges indefinitely would,
in a finite graph, eventually revisit a node, producing a cycle), and
repeatedly removing sources and appending them to the output list produces a
valid topological order.

This constructive procedure is, in fact, an algorithmic description of
exactly what a DAG-based scheduler does at runtime: a task with no
unsatisfied predecessor is, in DAGonStar's vocabulary, in the
\texttt{READY} state (\Cref{sec:status-model}); the engine's threaded
execution model effectively performs the source-removal procedure
\emph{concurrently}, starting every \texttt{READY} task as a thread, and
transitioning a waiting task to \texttt{READY} the moment its last
predecessor reaches \texttt{FINISHED}. \texttt{Validate\_WF()}'s cycle
rejection is therefore not an arbitrary safety check; it is the exact
precondition the underlying scheduling theorem requires for \emph{any}
valid execution order, sequential or concurrent, to exist at all.

\begin{practicebox}
This formal correspondence has a practical consequence worth internalizing:
because topological order is generally \emph{not unique} when a DAG contains
independent branches (as in \Cref{fig:basic-dag}, where \texttt{Berlin} and
\texttt{Nairobi} could legally appear in either order, or interleaved), a
DAGonStar workflow author should never write task logic that silently
depends on one sibling task happening to run before another, unless that
ordering is encoded as an explicit or inferred dependency. If two tasks are
not connected by an edge, the engine offers no guarantee about their
relative execution order or interleaving.
\end{practicebox}

\section{Granularity: How Big Should a Task Be?}
\label{sec:task-granularity}

A question every new workflow author eventually confronts, and one the
project's documentation does not answer prescriptively, is how finely a
computation should be decomposed into tasks. This is a genuine engineering
trade-off, and it is worth reasoning about explicitly rather than leaving it
to instinct.

\textbf{Coarse-grained tasks} --- for instance, a single \texttt{BATCH} task
that runs an entire multi-stage shell pipeline internally --- minimize the
number of scratch directories, launcher scripts, and \texttt{workflow://}
references the author must manage, and minimize the staging overhead
incurred at each dependency edge. Their cost is a loss of the very
properties \Cref{ch:introduction} opened by praising: traceability and
recoverability degrade, because a failure partway through a coarse task
provides no checkpoint boundary (\Cref{ch:lessons-05-07}) and no
intermediate, individually inspectable scratch directory.

\textbf{Fine-grained tasks} --- one task per logical step, connected by
explicit \texttt{workflow://} edges --- maximize traceability (every
intermediate product is a named, inspectable file in its own scratch
directory), maximize recoverability (checkpointing can skip any individual
completed step), and maximize the parallelism the engine can exploit
(\Cref{sec:topological-ordering}). Their cost is increased staging overhead
--- every edge triggers a data-movement operation
(\Cref{sec:data-staging-architecture}) --- and a larger number of artifacts
for a human reader to track.

\begin{practicebox}
A reasonable default, consistent with the five-stage
prepare/run/validate/postprocess/publish pattern this book returns to
repeatedly (\Cref{ch:external-software}), is to draw task boundaries at
points where a scientist would naturally want to \emph{inspect} an
intermediate product, or at points where the engine's checkpoint or
parallelism guarantees provide a concrete operational benefit --- not at the
boundary of every individual shell command. A model run that internally
invokes three tightly coupled binaries with no scientifically meaningful
intermediate product is usually one task; a model run followed by an
independently useful post-processing step is usually two.
\end{practicebox}

\section{Chapter Summary}

This chapter established the conceptual vocabulary used throughout the rest
of the book: a workflow is a DAG of tasks; dependencies may be declared
explicitly or inferred from data references; the engine materializes every
task into an isolated scratch directory; and a small, uniform set of task
types --- spanning local shells, HPC schedulers, containers, clouds, web
requests, native Python functions, and AI completions --- share one
dependency and staging mechanism. It further situated this design within the
broader history of computational pipeline tools (shell scripts, build
systems, workflow engines), gave a formal grounding for why acyclicity
guarantees an executable schedule, and offered a first discussion of task
granularity as a recurring design decision. \Cref{ch:architecture} now turns
this vocabulary into a precise account of the implementation that realizes
it.

\begin{exercise}
Using the source-removal construction of \Cref{sec:topological-ordering},
manually compute \emph{all} valid topological orders of the DAG in
\Cref{fig:basic-dag}. How many are there, and what does this number tell you
about the maximum parallelism the engine could exploit if it had unlimited
threads available?
\end{exercise}

\begin{exercise}
Using only the vocabulary of this chapter, classify the following dependency
as task-oriented or data-oriented, and justify your answer: ``the
post-processing script must not run until the simulation has written its
restart file.''
\end{exercise}
